% Proposed isolated A.92. Frozen cycle30 sources remain unchanged.
\statement{A.92}{Proved}{Exact common-word and swap-duality images at every odd Alexander head}
Let \(F=\widehat F(x,y)\), let \(z=(xy)^{-1}\), and let
\(\iota\) interchange x and y. For an odd positive integer n put
\(A_n=\mathbb F_2[C_n]\), with specified generator T, and let
\(I_n\) be its augmentation ideal. The three actual cyclic-cover
maps, written with their additive module coordinate before their head, are
\[
\begin{array}{c|cc}
 &x&y\\ \hline
 \kappa_x&(1,1)&(0,T)\\
 \kappa_y&(0,T)&(1,1)\\
 \kappa_z&(0,T)&(T^{-1},T^{-1}).
\end{array}
\tag{A.92.1}
\]
Their targets are \(A_n\rtimes C_n\), with
\((P,T^a)(Q,T^b)=(P+T^aQ,T^{a+b})\). In the third map
z goes to \((1,1)\); the retained generators are respectively
y, x, and x. Restricting to the closed commutator subgroup gives
a continuous homomorphism
\[
 \pi_n:F'\longrightarrow I_n^3.
\]
It is surjective. Moreover, for every \(U,W\in I_n\) satisfying
\(\overline W=TW\), there is one actual \(u\in F'\) such that
\[
 \iota(u)=u^{-1},\qquad
 \pi_n(u)=(U,U,W),\qquad
 [u(x,y)]_x=U,\quad [u(y,z)]_x=W.
\tag{A.92.2}
\]
Here the bar sends T to \(T^{-1}\). Both conclusions hold
simultaneously for all odd n: a compatible triple, or a compatible
pair in (A.92.2), is realized by a single full profinite source word.
These are source-word realization statements. They do not assert
that the constructed word is canonical for an involution, satisfies
the pentagon, or has a specified Grothendieck--Teichm\"uller multiplier.

\begin{proof}
\emph{The exact image of the common word.}\par
Each map in (A.92.1) exists by the universal property of F.
The target commutator subgroup lies in its augmentation ideal,
so the three coordinates of \(F'\) lie in \(I_n\). For
\(c=[x,y]=xyx^{-1}y^{-1}\), direct multiplication gives
\[
       \pi_n(c)=(1-T,T-1,1-T^{-1}).
\tag{A.92.3}
\]
Let \(R_n=\mathbb F_2[C_n\times C_n]\), with generators
X,Y recording conjugation by the original x,y. A conjugate
\(x^ay^b c y^{-b}x^{-a}\) multiplies the three coordinates
in (A.92.3) by \(T^b,T^a,T^{a-b}\). The closed normal
closure of c is exactly \(F'\): imposing \([x,y]=1\)
produces the universal two-generated profinite abelian group.
As the three coordinate groups are abelian, it follows that
\[
 \pi_n(F')=
 \left\{
 \begin{pmatrix}
 (1-T)P(1,T)\\ (T-1)P(T,1)\\
 (1-T^{-1})P(T,T^{-1})
 \end{pmatrix}:P\in R_n\right\}.
\tag{A.92.4}
\]
Every monomial of P is realized by the specified conjugate of c;
their finite product gives its sum. Conversely normal generation
gives no additional coordinates, so (A.92.4) is equality.

For a generator X of a cyclic order-n factor write
\(N_X=1+X+\cdots+X^{n-1}\), and similarly for Y, XY, and T.
Since n is odd, each norm is an idempotent with augmentation one.
Comparison of adjacent coefficients proves
\[
 \ker(T-1:A_n\to A_n)=\mathbb F_2N_T.
\]
Its intersection with \(I_n\) is zero. Thus multiplication
by \(T-1\), and also by \(1-T^{-1}\), is an automorphism
of the finite space \(I_n\).

The following three explicit elements separate the three lines
used in (A.92.4):
\[
 E_x=N_X(1-N_Y),\qquad
 E_y=N_Y(1-N_X),\qquad
 E_z=N_{XY}(1-N_X).
\tag{A.92.5}
\]
Their evaluations at \((1,T),(T,1),(T,T^{-1})\) are
\[
\begin{array}{c|ccc}
 & (1,T)&(T,1)&(T,T^{-1})\\ \hline
 E_x&1-N_T&0&0\\
 E_y&0&1-N_T&0\\
 E_z&0&0&1-N_T.
\end{array}
\tag{A.92.6}
\]
For example the third evaluation of \(N_{XY}\) is n, hence
one, while \(N_T^2=N_T\). On \(I_n\), the operator
\(1-N_T\) is the identity. For a prescribed triple
\((A_x,A_y,A_z)\in I_n^3\), take the unique \(Q_x,Q_y,Q_z\in I_n\)
with
\[
 (1-T)Q_x=A_x,\qquad (T-1)Q_y=A_y,
 \qquad (1-T^{-1})Q_z=A_z.
\]
Then
\[
       P=E_xQ_x(Y)+E_yQ_y(X)+E_zQ_z(X)
\]
realizes the prescribed triple in (A.92.4). This proves
surjectivity for every odd n, including the vacuous n=1 case.

\emph{Exact duality, rather than separate module words.}\par
On the three coordinates the actual swap acts as
\[
 \iota:(A_x,A_y,A_z)\longmapsto
                   (A_y,A_x,T^{-1}\overline{A_z}).
\tag{A.92.7}
\]
The last factor follows from the literal identity
\(\iota(z)=x^{-1}zx\) and inversion of the retained x-head;
it is not an unweighted interchange. Put
\(j(P)=T^{-1}\overline P\). On the coefficient basis it
sends index r to \(-r-1\) modulo n. There is exactly one
fixed index, \(r_0=(n-1)/2\), and all other orbits have size
two. A j-invariant vector in \(I_n\) has zero coefficient
at that fixed index, because all its paired coefficients have
even total augmentation. It is therefore a sum of paired norms.
A norm preimage can be chosen in \(I_n\): if its augmentation
is one, add \(T^{r_0}\), which has augmentation one and norm
zero. Hence
\[
                         (1+j)I_n=I_n^j.
\tag{A.92.8}
\]

Given the U,W in the statement, choose \(P_z\in I_n\) with
\(P_z+j(P_z)=W\). Surjectivity supplies an actual
\(w\in F'\) with \(\pi_n(w)=(U,0,P_z)\). Set
\[
                              u=w\iota(w)^{-1}.
\]
This word satisfies \(\iota(u)=u^{-1}\) in the full source,
and (A.92.7) gives \(\pi_n(u)=(U,U,W)\).
Finally \(\kappa_x(y)=(0,T)\) and
\(\kappa_x(z)=(T^{-1},T^{-1})\), exactly the two
\(\kappa_z\)-images of x,y. Thus the third coordinate of
u is literally the first coordinate of \(u(y,z)\).
All the based substitutions in (A.92.2) are thereby retained.

\emph{One full word for all odd heads.}\par
The maps are compatible under divisibility of n. Put
\[
 A_{\rm odd}=\varprojlim_{n\ \mathrm{odd}}A_n
   =\mathbb F_2\Bigl[\Bigl[\prod_{p\ \mathrm{odd}}\mathbb Z_p\Bigr]\Bigr],
 \qquad I_{\rm odd}=\ker(\operatorname{aug}).
\]
The preimages of any compatible triple are nonempty closed
subsets of compact \(F'\). Any finite collection is controlled
by its odd least common multiple, so the finite-intersection
property produces one full word realizing them all.
For the duality assertion take w as the compact variable and
impose the finite-coordinate conditions on
\(w\iota(w)^{-1}\). Their nonemptiness follows from the
construction above, and the same least-common-multiple argument
applies. No independent compatible selection of norm preimages
is needed.

The theorem concerns the same full source word in all three
modules. It proves that these common-word and exact-duality
constraints alone add no odd-head linear condition beyond the
stated augmentation and weighted invariance. Actual canonical
integrability, the full marked pentagon, and the nonlinear
\(\Xi_d/\Omega_\alpha\) equation remain separate; none is
inferred from this realization. A.90's criterion for actual
full-GT data is unchanged.
\end{proof}
